\chapter*{ЗАКЛЮЧЕНИЕ}
\addcontentsline{toc}{chapter}{ЗАКЛЮЧЕНИЕ}

В рамках выполнения курсовой работы была успешно спроектирована
и реализована база данных для хранения и анализа спортивной статистики.
Поставленная цель достигнута в полном объёме посредством решения
сформулированных задач:
\begin{enumerate}
    \item выполнен сравнительный анализ предметной области, сформулированы требования
          и ограничения к системе. Определен перечень метрик и формализована
          математическая методология их расчета. Выделены ключевые сущности и роли,
          построены диаграммы вариантов использования и диаграмма сущность-связь, выбрана
          реляционная модель данных;
    \item спроектирована структура базы данных в третьей нормальной форме. Разработаны
          алгоритмы автоматизированной обработки и расчёта аналитических показателей.
          Спроектированы ролевая модель доступа на уровне СУБД, серверное программное
          обеспечение для взаимодействия с БД и подсистема кэширования
          с префиксной инвалидацией ключей;
    \item обоснован выбор средств реализации, реализована структура базы данных
          средствами DDL и серверное приложение, предоставляющее API из двадцати двух
          маршрутов. Подготовлен набор тестовых данных объемом более ста пятидесяти
          тысяч фактов. Выполнено тестирование системы, доказавшее корректность
          работы и математическую точность расчёта показателей;
    \item проведено исследование влияния объема данных, стратегий индексации
          и уровня нагрузки на производительность выполнения аналитических запросов.
          Доказано, что применение целевых одиночных индексов сокращает среднее время
          выполнения аналитических выборок
          на~50\% и радикально снижает 95-й перцентиль задержки. Обращение
          к предвычисленной витрине ускоряет получение данных примерно в~31{,}5~раза
          по сравнению с прямой агрегацией, а использование in-memory кэширования
          снижает медианное время ответа сервера в~3{,}9~раза. На основе полученных
          результатов сформулированы рекомендации по оптимизации и масштабированию
          системы при росте конкурентной нагрузки.
\end{enumerate}

Разработанный программный комплекс представляет собой готовое
масштабируемое решение, применимое для обеспечения деятельности
аналитических отделов спортивных организаций.
